\section{Discussion}\label{sec:discussion}

The reference pipeline described in this article is not, in the first instance,
a verification technique; it is a proposal to reorganise the software
development \emph{process} around a machine-checkable specification. The
technical mechanisms---inference, elaboration, the synthesis--verification
loop, compositional propagation---are in service of a process change, and it is
that change, rather than any single algorithm, that constitutes the
article's central contribution to the software development process. This section therefore steps back from
the architecture and asks what the lifecycle looks like once a formal contract
layer sits at the pivot between human intent and machine output: what the unit
of work becomes, where trust is anchored, how dependency evolution is handled,
how the approach is adopted incrementally, and---honestly---where it will and
will not repay the effort.

\subsection{The process shift: from reviewing code to reviewing
specifications}\label{subsec:disc-processshift}

The most consequential implication is a relocation of the human review
activity. In the prevailing process, the durable record of intent is the code
itself, and quality control is exercised by reading code in review. When an
agent can regenerate an implementation cheaply, the implementation becomes a
disposable build output and the scarce, durable artefact is the contract it was
built and verified against. Review then migrates up the stack: the reviewer
inspects the \emph{specification delta}, not the diff of generated statements.
This is a qualitatively easier object to reason about, because a contract states
\emph{what} must hold over all inputs rather than \emph{how} one particular
realisation computes it, and because the implementation has already been proved
to satisfy whatever the contract says by OpenJML's extended static
checking~\cite{cok2014openjml}, covering all inputs rather than the sampled
inputs a test suite exercises.

Correspondingly, the unit of change becomes a \emph{spec-delta plus a discharged
proof obligation} rather than a code patch. A pull request in this model carries
the changed obligations, their provenance and assurance tags, and the verifier's
verdict that the new implementation satisfies them. The reviewer adjudicates the
contract change---is this the behaviour we want?---and delegates the question of
whether the code conforms to the verifier. This is the sense in which trust
becomes \emph{transferable}. In conventional review, trust is established by a
qualified human eyeballing the implementation, an act that does not compose and
does not transfer: a reviewer's approval is an opinion, not a checkable artefact.
When conformance is established by a discharged proof obligation, the trust is
attached to the artefact and can be re-checked by anyone, including a machine.
It is precisely this property that would let an agent merge its own work without
a human in the conformance loop: not because the agent is trusted, but because
its output carries an independently re-verifiable certificate that it satisfies a
human-ratified contract. The human authority is exercised once, over the
contract; it is not re-spent on every implementation.

Two caveats keep this from becoming a slogan. First, transferable trust is only
as strong as the contract: a discharged proof against a weak or wrong contract
certifies nothing useful, which is why the independent-authority constraint and
behaviour-level ratification described in the elaboration layer are load-bearing
rather than decorative. Second, the verdict is honest only if it carries its
assurance level. A merge gate must distinguish an obligation that was
\texttt{proved} from one that was merely \texttt{bounded-checked} or
\texttt{tested} after the proof timed out; the tiered-assurance tagging exists so
that ``verified'' on a pull request is an auditable claim rather than a
reassuring word.

\subsection{Dependency evolution as a spec-diff event}\label{subsec:disc-depupgrade}

The second process change follows from compositionality. Because contracts
compose, a change to a callee's contract mechanically recomputes the proof
obligations of its callers, and the verifier reports which callers no longer
discharge---without executing any of them. Internal refactoring and external
dependency upgrades are then the same operation viewed at different scopes: in
both cases a callee specification moves and the blast radius is the set of call
sites whose obligations break. This recasts the perennial and poorly served
question ``which of my call sites does this library upgrade break
behaviourally?'' as a first-class, answerable lifecycle event. The existing
literature on breaking changes is overwhelmingly signature-level---tools such as
APIDiff and the Maven-Central replication studies detect that a method's
\emph{shape} changed~\cite{brito2018apidiff,ochoa2022breaking}---while the harder
and more common problem is the \emph{behavioural} breaking change, where the
signature is stable but the contract is not, and which the same line of work
identifies as the dominant unsolved
case~\cite{mostafa2017behavioral,jayasuriya2024understanding}. Semantic
versioning encodes a promise about compatibility but cannot check
it~\cite{preston2013semver,raemaekers2017semver}. A contract-diff over composed
specifications turns that promise into a verification condition: an upgrade is
behaviourally compatible for a given client exactly when the client's obligations
still discharge against the new callee contract. This is the lifecycle
expression of the thesis link, with the compositional refinement validated in
our prior work~\cite{edmonds_article3} promoted from an offline corpus
measurement to an online change-impact capability, and it is the point at which
the formal-methods-adoption and library-compatibility strands of the programme
meet in a single mechanism.

\subsection{An adoption path that earns trust before it demands it}\label{subsec:disc-adoption}

A reference architecture is worthless if it can only be entered greenfield, and
the overwhelming majority of software is brownfield. The adoption path is
therefore designed to be incremental along two axes. The first is
\emph{scope}: verification does not scale to an entire legacy application, so the
contract layer is introduced at a frontier---the module under active change, the
public API surface, and the modules referenced by live tickets---while the
remainder is treated as unspecified and governed only by weak inherited default
contracts. The boundary contracts at this frontier are exactly where
compositional propagation attaches, so the specified region can grow outward
along call edges rather than requiring a big-bang formalization. The second axis
is \emph{sequence}. The first pass over a brownfield codebase is
\emph{retroactive}: the startup phase recovers a draft contract from the code,
triangulates it against tickets and documentation, and surfaces the
disagreements as candidate bugs and stale docs to be resolved with a human;
only after this establishes a trusted baseline do \emph{prospective} passes
introduce new behaviour. The tool thus earns its keep on existing code---finding
real discrepancies in a system the team already owns---before it is permitted to
change anything. Inherited default contracts carry much of the adoption burden:
organisation-level policies (public methods reject null unless stated, monetary
values are never negative, returned collections are never null) are specified
once and inherited everywhere, so that the per-method ratification labor is
spent only on departures from default. The human role, across all of this, moves
up the stack: away from writing and reviewing implementations and toward
adjudicating contracts, ratifying behaviour through concrete examples, and
arbitrating the disagreements the triangulation surfaces.

\subsection{What is genuinely novel}\label{subsec:disc-novelty}

It is important to be precise about what distinguishes this proposal from the
substantial bodies of work it builds on, because individually most of its
ingredients are not new. Counterexample-guided synthesis is decades
old~\cite{solarlezama2008sketching,alur2013sygus}; closed-loop verified code
generation with a language model and a deductive verifier has recent and capable
instances~\cite{sun2024clover,misu2024dafny,mugnier2025dafnypro}; specification
inference, both dynamic and learned, is a mature
field~\cite{ernst2007daikon,ma2025specgen,endres2024nl2postcond}; and
spec-driven agentic lifecycles are emerging in
industry~\cite{githubspeckit,awsaidlc,awskiro}. The novelty is not in any one of
these but in three commitments that the emerging agentic lifecycles in
particular do not make. The first is insistence on a \emph{machine-checkable}
contract at the pivot, in place of the natural-language specification those
lifecycles treat as ground truth; this is what closes the otherwise open
generate-then-eyeball loop with a machine gate that ranges over all inputs. The
second is the \emph{independent-authority constraint}: the contract a code is
verified against must not be authored by the agent that wrote the code,
dissolving the circularity that afflicts self-grading approaches where the same
model produces an implementation and the tests that judge it---a hazard the
test-oracle literature has begun to flag explicitly for LLM-generated
oracles~\cite{liu2026oracle,barr2015oracle}. The third is the elaboration layer's
treatment of the \emph{ratification gap}: rather than asking a human to confirm
formal text they used natural language precisely to avoid, the contract is
ratified at the level of behaviour through pinned, human-confirmed examples, and
its strength is policed by an adversarial elaborator that hunts for scenarios
consistent with the stated intent yet permitted by an over-weak contract. The
combination---a verification spine, an authority separation, and a
behaviour-level human gate---is, to our knowledge, not assembled in prior agentic
or spec-driven work, and it is that assembly, instantiated on validated
components, that constitutes the contribution.

\subsection{Where the approach will and will not pay off}\label{subsec:disc-payoff}

Candor about applicability is owed to the reader, both because the claim of this
article is feasibility rather than measured effect and because the failure modes
are predictable. The approach pays off where behaviour is precisely statable and
mechanically checkable: data-structure and library code, numeric and bounds
contracts, null- and emptiness-discipline, and the loop patterns---maximum,
minimum, linear search, summation---for which the inference instrument already
produces discharging invariants. These are the method classes that verify green
today and the natural first targets for instantiation. It pays off least where
intent resists formalization (subjective or aesthetic requirements, distributed
and timing-dependent behaviour, performance properties outside a functional
contract) and where the underlying decision procedures do not terminate; the
Z3-bound timeouts on multiplicative reasoning, array theory, and rich
preconditions observed across our verification corpus are real, known, and not
papered over---they are precisely the cases the tiered-assurance mechanism exists
to degrade gracefully rather than to deny.

Several limitations bear directly on the process claims above and must be stated
plainly. Inferred contracts skew weak, describing what code does rather than
what it should, so ratification is genuine human labor; the mitigation is
inherited defaults plus API-surface prioritisation, not an assertion that the
labor disappears. Inferred contracts are sometimes wrong or over-strong, which
manifests as false candidate bugs; this is why every recovered clause is offered
as a \emph{candidate} and why human ratification, not the inference, owns the
truth. Translating ticket and documentation intent into formal clauses without a
human is imperfect, so those clauses remain unconfirmed and drive the dialogue
rather than gating it, and linking tickets and docs to code is noisy enough that
it is credible only over a curated module rather than a whole repository. Above
all, a wrong contract verified against is worse than no contract, because it
launders an error into a proof; this single observation is why the independent
authority and behavioural ratification are not optional refinements but the load
bearing structure of the whole proposal. Finally, what verification establishes
over a brownfield codebase must not be overstated: proving code against a
code-derived contract is characterisation---a faithful-description baseline and a
toolchain sanity check---and is nearly circular; it does not prove the absence of
bugs. Correctness relative to intent emerges only from the ticket and
documentation streams and the human ratification on top of them, and bugs appear
as \emph{disagreements} between what the code does and what those streams say it
should. The full empirical adjudication of these claims on external production
codebases, including an industrial action-research case study under appropriate
ethics review, is future work; this article claims that the architecture is
realistic and instantiable on validated components, not that its benefits have
yet been measured at scale.
